\documentclass[pdflatex,sn-mathphys-ay]{sn-jnl}

\usepackage[utf8]{inputenc}%
\usepackage[T1,T2A]{fontenc}%
\usepackage[english,russian]{babel}%
\usepackage{graphicx}%
\usepackage{multirow}%
\usepackage{amsmath,amssymb,amsfonts}%
\usepackage{amsthm}%
\usepackage{mathrsfs}%
\usepackage[title]{appendix}%
\usepackage{xcolor}%
\usepackage{textcomp}%
\usepackage{manyfoot}%
\usepackage{booktabs}%
\usepackage{array}%
\usepackage{longtable}%
\usepackage{calc}%
\usepackage{newunicodechar}%
\providecommand{\tightlist}{\setlength{\itemsep}{0pt}\setlength{\parskip}{0pt}}
\newunicodechar{§}{\S}
\newunicodechar{→}{\ensuremath{\rightarrow}}
\newunicodechar{↔}{\ensuremath{\leftrightarrow}}
\newunicodechar{×}{\ensuremath{\times}}
\newunicodechar{·}{\ensuremath{\cdot}}
\newunicodechar{⊃}{\ensuremath{\supset}}
\newunicodechar{°}{\ensuremath{^{\circ}}}
\newunicodechar{«}{\guillemotleft}
\newunicodechar{»}{\guillemotright}

\raggedbottom

\begin{document}

\title[Дополнительные материалы]{Supplementary Materials for «Landauer Efficiency of Self-Modeling: An Operational Scale of Vitality»}

\author*[1]{\fnm{Александр} \sur{Андриишин}}\email{alexander@andriishin.com}

\affil*[1]{\orgname{Independent Researcher}, \orgaddress{\city{Москва}, \country{Россия}}}

\maketitle

\section{\texorpdfstring{S2.1 Операционализация \(I_{\text{pred}}\): полный протокол}{S2.1 Операционализация I\_\{\textbackslash text\{pred\}\}: полный протокол}}\label{s2.1-ux43eux43fux435ux440ux430ux446ux438ux43eux43dux430ux43bux438ux437ux430ux446ux438ux44f-i_textpred-ux43fux43eux43bux43dux44bux439-ux43fux440ux43eux442ux43eux43aux43eux43b}



Прямое измерение \(I_{\text{pred}}\) в смысле формулы (1a) основного текста требует доступа ко всем внутренним степеням свободы системы и для биологических, городских и космических систем технически невозможно. Стандартизованная операционализация фиксирует для каждого класса систем четыре элемента.



(а) \emph{Сенсорный канал} \(x_t\) --- наблюдаемые входные состояния. Для бактерии --- концентрация лиганда у хемотаксических рецепторов; для города --- поток ресурсов через инфраструктуру (энергия, материя, информация); для биосферы --- поток солнечной радиации и геохимические потоки.



(б) \emph{Внутренний канал} \(s_t\) --- измеримые сводки конфигурации \(M_t\). Для бактерии --- уровень метилирования адаптационного механизма; для города --- административно-инфраструктурные параметры (плотность коммуникаций, плотность решений, профиль бюджета); для биосферы --- атмосферный состав и пулы биомассы.



(в) \emph{Выбор окон} прошлого/будущего с указанием \(\tau\). Для бактерии --- характерное время реакции, порядка нескольких поколений; для города --- квартал или год; для биосферы --- геологический эон.



(г) \emph{Метод оценки взаимной информации.} Прямой расчёт \(I(s_t; x_{t+\tau})\) требует знания совместного распределения \(p(s, x)\), недоступного для биологических, городских и космических систем. Используются оценки по конечной выборке; выбор оценщика зависит от природы данных (дискретные/непрерывные, низко-/высокоразмерные, объём выборки \(N\)). Три класса методов, применяемые в настоящей работе:



\emph{Plug-in оценка с коррекцией смещения.} Эмпирические частоты \(\hat p(s, x) = N(s, x)/N\) подставляются в формулу \(\hat I = \sum_{s, x} \hat p(s, x) \log[\hat p(s, x)/(\hat p(s) \hat p(x))]\). Главная проблема --- систематическая переоценка MI при малых выборках (\emph{finite-sample bias}) порядка \((m-1)/(2N \ln 2)\) бит, где \(m\) --- число занятых ячеек в гистограмме. Стандартные коррекции: \textbf{Miller--Madow} (Miller 1955) --- прибавление к наивной эмпирической оценке энтропии члена \((m-1)/(2N \ln 2)\) для устранения систематического занижения, корректно при \(m \ll N\); \textbf{Grassberger--Schurmann} (Grassberger 2003) --- оценка энтропии через digamma-функцию, \(\hat H_{\text{GS}} = \log N - \frac{1}{N} \sum_s n_s \psi(n_s)\), дающая меньший bias при средних значениях \(m/N\). Plug-in оптимальна для \textbf{дискретных биологических данных} низкой/средней размерности (статусы метилирования рецепторов, дискретные клеточные состояния, \(\dim \le 10\), выборка \(N \sim 10^3\)--\(10^5\)): естественно дискретная природа сигнала, отсутствие проблем с разбиением на интервалы.



\emph{Kraskov--Stögbauer--Grassberger (KSG; Kraskov et al.~2004).} \(k\)-nearest-neighbour оценка для \textbf{непрерывных данных}, не требующая разбиения на интервалы. Формула: \(\hat I_{\text{KSG}} = \psi(k) + \psi(N) - \langle \psi(n_x + 1) + \psi(n_y + 1) \rangle\), где \(\psi\) --- digamma-функция, \(n_x\) и \(n_y\) --- число соседей в маргинальных \(L_\infty\)-окрестностях, размер которых задаётся \(k\)-м ближайшим соседом в совместном пространстве. Адаптивно учитывает локальную плотность распределения. Параметр \(k\) контролирует \emph{bias--variance trade-off}: меньшие \(k\) → меньше bias, больше дисперсии; рекомендация авторов \(k = 4\). Применима для \textbf{непрерывных временных рядов средней размерности} (\(\dim \le 50\)), выборка \(N \sim 10^3\)--\(10^6\) --- стандарт для нейрофизиологических, физиологических и эконометрических временных рядов.



\emph{Mutual Information Neural Estimation (MINE; Belghazi et al.~2018).} Использует \textbf{Donsker--Varadhan дуальное представление} KL-дивергенции: \(I(X; Y) = \sup_T \{\mathbb{E}_{p_{XY}}[T(x, y)] - \log \mathbb{E}_{p_X \otimes p_Y}[\exp T(x, y)]\}\), где супремум берётся по всем измеримым функциям \(T: \mathcal{X} \times \mathcal{Y} \to \mathbb{R}\). Функция \(T\) («критик») параметризуется нейросетью и обучается стохастическим градиентным спуском по mini-batch с тождеством, что MI равна нижней границе по любому \(T\). Применима к \textbf{высокоразмерным непрерывным данным} (изображения, LLM-эмбеддинги, размерность \(> 10^2\)), где KSG страдает от curse of dimensionality (объём \(L_\infty\)-окрестности растёт как \(r^d\), и оценка плотности через kNN деградирует). Цена --- отсутствие гарантий точной сходимости при недоучённом критике и существенные вычислительные затраты (часы--сутки на GPU для типичной выборки \(N \sim 10^5\), \(\dim \sim 10^3\)).



\textbf{Метод по умолчанию для каждого класса данных рекомендуется следующий.}



\begin{center}\small
\begin{tabular}{@{} >{\raggedright\arraybackslash}p{0.3333\linewidth} >{\raggedright\arraybackslash}p{0.3333\linewidth} >{\raggedright\arraybackslash}p{0.3333\linewidth}@{}}
\toprule

\begin{minipage}[b]{\linewidth}\raggedright

Класс данных

\end{minipage} & \begin{minipage}[b]{\linewidth}\raggedright

Метод

\end{minipage} & \begin{minipage}[b]{\linewidth}\raggedright

Параметры по умолчанию

\end{minipage} \\

\midrule

Дискретные биологические (\(\dim \le 10\)) & Plug-in с Miller--Madow коррекцией & sanity check через KSG с \(k = 4\) \\

Непрерывные временные ряды (\(\dim \le 50\)) & KSG & \(k = 4\) \\

Высокоразмерные непрерывные (\(\dim > 10^2\)) & MINE & критик --- MLP с 3 скрытыми слоями по 512 нейронов, обучение \(\ge 10^4\) batch-update \\

\bottomrule
\end{tabular}
\end{center}


\textbf{Диагностика чувствительности.} При расхождении более 30\% между двумя методами на одной выборке --- варьировать параметры, специфичные для оценщика:



\begin{itemize}

\tightlist

\item

  \emph{KSG:} варьировать \(k \in \{3, 5, 8, 10\}\) и проверять \emph{плато} --- стабильное значение оценки в среднем диапазоне \(k\). Отсутствие плато (монотонный дрейф оценки по \(k\)) указывает на недостаточный объём выборки или сильную локальную неоднородность плотности.

\item

  \emph{Plug-in:} варьировать размер бина при разбиении на интервалы непрерывных данных, сравнивать коррекции Miller--Madow и Grassberger--Schurmann. Расхождение между коррекциями \(> 10\%\) --- сигнал, что \(m/N\) слишком велик и нужен KSG или MINE.

\item

  \emph{MINE:} варьировать архитектуру критика (число слоёв, ширина), проверять \emph{сходимость bound} по эпохам обучения (стабильное значение в последних \(\sim 10^3\) batch-updates), регуляризацию (gradient clipping против разрастания градиентов в exp-члене). Не-сходящийся bound, отрицательные оценки или сильная зависимость от random seed --- сигнал к смене архитектуры или к переходу на KSG для нижне-размерных подпроекций.

\end{itemize}



Для систем с активным sampling (бактерия с моторной обратной связью на градиент, город с административной интервенцией) операциональная MI-оценка проводится при фиксированной политике \(\pi\): \(I_{\text{pred}}(\pi) = I(s_t; x_{t+\tau} \mid \pi)\). Без conditioning DPI-неравенство \(I(s_t; x_{t+\tau}) \le I(x_{t-\tau:t}; x_{t:t+\tau})\) не гарантировано, и MI-оценка смешает predictive information с эффектом активного управления (см. main § 2.1 DPI-условие).



Прокси не суть определение \(I_{\text{pred}}\) --- они операционализация. Та же дисциплина действует в стандартной физической измерительной практике: температура измеряется не прямым подсчётом скоростей молекул, а через калиброванный прокси --- расширение ртути. Расхождение между разными прокси для одной системы --- сигнал к уточнению прокси, не к отказу от концепта.



\emph{Почему predictive, а не stored information.} Прецедент термодинамической связи predictive information с диссипацией задан в работе Still, Sivak, Bell, Crooks о термодинамике предсказания (Still et al.~2012). Там введён термин \textbf{nostalgia} для бесполезно удерживаемой системой прошлой информации, не имеющей предсказательной силы: \(I_{\text{nostalgia}} := I_{\text{stored}} - I_{\text{pred}} \ge 0\) по DPI; равенство достигается, когда внутренняя статистика \(M_t\) минимально достаточна для предсказания (оптимум Information Bottleneck (Tishby et al.~1999)). Still et al.~показывают, что \(I_{\text{nostalgia}}\) задаёт нижнюю границу на полную диссипацию системы. В нестационарной среде разрыв \(I_{\text{nostalgia}} > 0\) типичен и связан с избыточной диссипацией; в стационарной с фиксированным \(M < \infty\) разрыв сохраняется, но ограничен сверху ёмкостью памяти (Still et al.~2012). Настоящая рамка наследует прецедент Still'а и добавляет ключевое требование: и числитель, и знаменатель \(\eta_L\) относятся к одной и той же физической системе --- модель и диссипация принадлежат \(X\). Это требование self-payment отделяет \(\eta_L\) от классической постановки термодинамики предсказания, где модель и диссипатор могут быть разными контурами.



\subsection{Лемма о стационарном стирании: полное доказательство}\label{ux43bux435ux43cux43cux430-ux43e-ux441ux442ux430ux446ux438ux43eux43dux430ux440ux43dux43eux43c-ux441ux442ux438ux440ux430ux43dux438ux438-ux43fux43eux43bux43dux43eux435-ux434ux43eux43aux430ux437ux430ux442ux435ux43bux44cux441ux442ux432ux43e}



Main § 2.1 содержит informal-summary Леммы и краткое формальное утверждение с двумя частями (поток-граница и сток-граница) плюс плотный мета-параграф о соотношении доказанного и постулируемого. Здесь приводится полный набросок доказательства с детальным разбором алгоритмической части (шаги (i)--(iii), (v)) и обсуждением постулата стационарного учёта (шаг (iv)).



\textbf{Условия Леммы.} \(X\) в стационарном эргодическом режиме удерживает predictive information \(I_{\text{pred}}\) на окне \(\tau\); память системы конечна (\(M < \infty\) битов); поток новой predictive information строго положителен (\(\dot{I}_{\text{pred}} > 0\)). Лемма содержит две логически независимые границы.



\textbf{Часть (а): поток-граница (справочная).} Темп предсказательного обновления ограничен сверху скоростью энтропии источника: \(\dot I_{\text{pred}} \le h_\mu\), где \(h_\mu\) --- скорость энтропии источника. Прямое следствие data processing inequality при \(M_t = f(X_E^{(-\infty, t]})\) как функции прошлого сенсорного канала и стационарности (с conditioning на политику в случае активного sampling --- см. ниже). Равенство достигается в пределе оптимального предсказания, когда вся энтропия источника предсказательно релевантна. I.i.d.-источник (independent and identically distributed --- последовательность независимых одинаково распределённых отсчётов, в которой прошлое не несёт никакой информации о будущем) реализует противоположный предел: \(h_\mu > 0\) при \(I_{\text{pred}} = 0\) и \(\dot{I}_{\text{pred}} = 0\). Граница (а) --- справочная верхняя оценка на \(\dot I_{\text{pred}}\) для связи с классической линией Биалека--Неменмана--Тишби (Bialek et al.~2001); \textbf{в выводе нижней границы диссипации она не используется} и приводится лишь для контекста.



\textbf{Часть (б): сток-граница (центральная).} При \(M < \infty\) и положительности энтропийного потока через систему за окно \(\tau_d\) --- при истинности постулата стационарного учёта (шаг (iv) ниже) --- обязательное стирание \(\ge I_{\text{pred}}\) битов и диссипация \(\ge I_{\text{pred}} \cdot k_B T \ln 2\) свободной энергии. Сток-граница условна. Алгоритмическая часть (шаги (i)--(iii)) доказывает, что память должна освобождать \(\dot I_{\text{pred}}\) битов в единицу времени. Переход от этого освобождения к ландауэровской плате (шаг (iv)) принимается как постулат стационарного учёта в отсутствие глобального обратимого расвычисления (Bennett uncomputation) (Bennett 1973, 1982).



\emph{Набросок доказательства (б):}



\begin{enumerate}

\def\labelenumi{(\roman{enumi})}

\item

  Стационарность распределения.

\item

  При выборе окна \(\tau = \tau_d\) тождество \(\dot I_{\text{pred}} \cdot \tau_d = I_{\text{pred}}\) выполняется по определению \(\tau_d\) --- это переименование, а не содержательный шаг. Эргодичность обеспечивает существование стационарных значений обеих величин.

\item

  \emph{Алгоритмическая нижняя граница на содержимое памяти.} Связь шенноновской взаимной информации \(I_{\text{pred}}\) с алгоритмической сложностью \(M_t\) требует двух шагов. \textbf{Шаг 1} (Shannon-уровень): по неравенству \(I(X; Y) \le \min(H(X), H(Y))\), где \(H(\cdot)\) --- шенноновская энтропия дискретной случайной величины (\(H(X) = -\sum_x p(x) \log_2 p(x)\), в битах при логарифме по основанию 2; Shannon 1948), имеем \(H(M_t) \ge I(M_t; X_E^{[t, t+\tau_d]}) = I_{\text{pred}}\). \textbf{Шаг 2} (алгоритмическая нижняя граница на \(M_t\)): стандартная связь между средней колмогоровской сложностью и шенноновской энтропией источника (Cover and Thomas 2006; Li and Vitányi 2008) даёт \[\langle K(M_t) \rangle \ge H(M_t) + O(\log) \ge I_{\text{pred}} + O(\log)\] для типичных реализаций memory state, где \(\langle \cdot \rangle\) --- среднее по инвариантной мере \(\mu\) источника. Нижняя граница \(\langle K(M_t)\rangle \ge H(M_t) - O(1)\) --- это стандартное неравенство теории кодирования (ожидаемая колмогоровская сложность не ниже энтропии с точностью до аддитивной константы); второе неравенство --- Шаг 1 (DPI). Несущий вывод нижней границы памяти не использует равенство Brudno: он опирается лишь на \(H(M_t) \ge I_{\text{pred}}\) (Шаг 1) и на это стандартное кодовое неравенство. Нотационное пояснение: \(O(\log)\) --- поправка порядка \(\log \tau_d\) (член, ограниченный сверху \(C \cdot \log \tau_d\) для некоторой константы \(C\)), связанная с алгоритмическим описанием самого декомпрессора. Соответственно минимальное физическое представление \(M_t\) требует не менее \(I_{\text{pred}}\) битов памяти.

\end{enumerate}



\textbf{Мост Brudno (к шкале \(h_\mu\), часть (а)).} Отдельно от вывода нижней границы памяти, по теореме Brudno (1983) для эргодического источника с \(h_\mu > 0\) типичная траектория имеет колмогоровскую сложность \(K(\cdot)\): \[K(X_E^{[t-\tau_d, t]}) = h_\mu \cdot \tau_d + o(\tau_d) = H(X_E^{[t-\tau_d, t]}) + o(\tau_d)\] для \(\mu\)-почти всех траекторий. Нотационное пояснение: \(o(\tau_d)\) --- стандартное обозначение Ландау для члена с \(o(\tau_d)/\tau_d \to 0\) при \(\tau_d \to \infty\), то есть члена, асимптотически малого по сравнению с длиной окна. Это равенство связывает алгоритмическую сложность траектории \emph{источника} со скоростью энтропии \(h_\mu\) и служит мостом к классической шкале Биалека--Неменмана--Тишби (часть (а) Леммы), а не посылкой вывода нижней границы памяти \(M_t\). При \(M < \infty\) и стационарном притоке новой predictive информации с темпом \(\dot I_{\text{pred}} > 0\) ёмкость памяти насыщается за время \(\sim M/\dot I_{\text{pred}}\); дальнейшая работа канала требует освобождения \textbf{не менее} \(\dot I_{\text{pred}}\) битов predictive-релевантной информации в единицу времени (приходить может больше --- непредсказательный шум также подлежит стиранию, что только усиливает результат).



\begin{enumerate}

\def\labelenumi{(\roman{enumi})}

\setcounter{enumi}{3}

\item

  \emph{Постулат стационарного учёта.} Освобождение в (iii) реализуется как ландауэровское стирание \(\ge I_{\text{pred}}\) битов за окно \(\tau_d\) в любой реализации, не использующей глобальное обратимое расвычисление (Bennett 1973, 1982). Bennett-путь формально не запрещён. Time/space-tradeoff Bennett (1989) даёт: любое необратимое вычисление, требующее времени \(t\) (число логических шагов) и памяти \(S\) (битов), эмулируется обратимо с памятью \(S_R \sim S(1 + \log(t/S))\) при времени \(T_R \sim t^{1+\epsilon}/S^\epsilon\), где \(S_R\), \(T_R\) --- память и время обратимой симуляции, \(\epsilon > 0\) --- свободный сколь угодно малый параметр (меньшее \(\epsilon\) --- быстрее, но с большим префактором времени). Содержательно: логарифмический рост памяти симуляции, неограниченный при \(t \to \infty\) и фиксированном \(M\). Sagawa and Ueda (2010) и Parrondo et al.~(2015) формализуют обобщённое второе начало \(\langle\Sigma\rangle \ge -\langle\Delta I\rangle\) для feedback-controlled систем (Sagawa and Ueda 2009 --- частный случай measurement+erasure cost). Эта же рамка указывает на потенциальный контр-пример к границе: стирание памяти, \emph{коррелированной} со средой, допустимо ценой \(\langle W \rangle \ge k_B T(\Delta S - I_{\text{meas}})\), то есть ниже \(k_B T \ln 2\) за бит. Поскольку удерживаемая predictive-память коррелирована со средой по построению (\(I_{\text{pred}} = I(M_t; X_E) > 0\)), такой канал в принципе мог бы дать \(N_{\text{max}} < I_{\text{pred}}\) и \(\eta_L > 1\). Механизм, однако, требует контроллера, \emph{сначала} измеряющего коррелированный эталон с приобретением \(I_{\text{meas}}\); по симметрии измерение↔стирание (Sagawa and Ueda 2009) само измерение стоит \(\ge k_B T \cdot I_{\text{meas}}\), так что за полный стационарный цикл (приобретение корреляции + расходование её на дешёвое стирание) суммарная диссипация остаётся \(\ge k_B T \ln 2\) на бит. В петле самооплаты учёт внутренний: внешнего демона со свободным доступом к корреляции нет, а сама \(I_{\text{pred}}\) поддерживается тем же самооплачиваемым бюджетом, поэтому корреляцию в числителе нельзя одновременно потратить на удешевление знаменателя без двойного учёта (демон Максвелла внутри замкнутой системы). Постулат стационарного учёта формализует именно это замыкание цикла; условный статус границы при этом сохраняется. Полный учёт произвольного канала обновления \(M_t \to M_{t+1}\) при \(M < \infty\) и \(\dot I_{\text{pred}} > 0\) --- открытая техническая задача. В настоящей работе переход от (iii) к ландауэровской плате принимается как \textbf{постулат стационарного учёта}. Формулировка: при \(M < \infty\) и \(h_\mu > 0\), вне глобальной Bennett-асимптотики, канал обновления диссипирует не менее \(\dot I_{\text{pred}} \cdot k_B T \ln 2\) свободной энергии в единицу времени. Стираются старые элементы памяти --- типично \emph{nostalgia}-биты в смысле выше, не сами predictive биты. Граничный случай глобально-обратимой реализации отделён режимной оговоркой \emph{Bennett-режим и переопределение шкалы} в main § 2.1 и явно перенесён в условие (в) Утверждения об области определения.

\item

  Нижняя граница на одно стирание (Landauer 1961). При истинности постулата стационарного учёта (iv) суммарная диссипация на окне \(\tau_d\) составляет \(\ge I_{\text{pred}} \cdot k_B T \ln 2\), и \(\eta_L \le 1\).

\end{enumerate}



Для нестационарных систем с растущей памятью (биосфера, обучающаяся искусственная система) лемма требует учёта стоимости удержания старой информации, расширения памяти и поглощения новой; полная сумма с ландауэровской границей --- открытая техническая задача. Контекст более общих результатов --- Sagawa and Ueda (2010), Parrondo et al.~(2015), Still et al.~(2012).



\textbf{Расширенный статус Леммы.} Часть (а) Леммы и шаги (i)--(iii), (v) части (б) доказательны. Шаги (i)--(ii) --- тождества по определению \(\tau_d\) и стационарности. Шаг (iii) --- \emph{асимптотическое} следствие, опирающееся на (1) Shannon-уровневое неравенство \(H(M_t) \ge I(M_t; \cdot) = I_{\text{pred}}\) (DPI) и (2) стандартную нижнюю границу теории кодирования \(\langle K(M_t)\rangle \ge H(M_t) - O(1)\) (Cover and Thomas 2006; Li and Vitányi 2008) о том, что ожидаемая колмогоровская сложность не ниже энтропии. Равенство Brudno (1983) о совпадении средней колмогоровской сложности и скорости энтропии источника входит отдельно --- как мост к шкале \(h_\mu\) (часть (а) Леммы), не как посылка вывода нижней границы памяти. Соотношение \(\langle K(M_t)\rangle = H(M_t) + O(\log)\) выполнено для типичных реализаций memory state; формальный учёт допустимых режимов сжатия памяти выходит за рамки настоящей работы. Шаг (v) --- классический принцип Ландауэра. \textbf{Содержательно недоказанным остаётся шаг (iv)} --- переход от «память должна освободить \(\dot I_{\text{pred}}\) битов в единицу времени» к «это освобождение реализуется именно как ландауэровское стирание, не как обратимое расвычисление». Текущая литература формализует это для feedback-controlled систем (Sagawa and Ueda 2010; Parrondo et al.~2015) и для time/space-tradeoff обратимой симуляции (Bennett 1989), но не для произвольного канала обновления в стационарном режиме. В настоящей работе шаг (iv) принимается как \textbf{постулат стационарного учёта} и фиксируется в условии (в) Утверждения об области определения. Соответственно \(\eta_L \le 1\) --- условное утверждение: следствие принципа Ландауэра при истинности постулата (iv), а не безусловная теорема. Доказательство (iv) для произвольного канала обновления вынесено в раздел открытых задач (main § 5).



\textbf{Bennett-режим и переопределение шкалы: расширенное обсуждение.} Main § 2.1 содержит сжатое scope-statement о Bennett-режиме. Полная аргументация: для систем в Bennett-асимптотике обратимых вычислений (Bennett 1973, 1982) плата на бит уходит ниже Ландауэра, и \(N_{\text{max}}\) теряет интерпретативный смысл. При \(E_{\text{actual}} \to 0\) и конечном \(I_{\text{pred}}\) отношение \(\eta_L\) формально расходится. В Bennett-режиме шкала переопределяется через число логических, а не стираемых битов --- операционально другая величина, требующая отдельной калибровки. В настоящей работе мы остаёмся в необратимом режиме как стандартной идеализации биологических и инженерных систем.



Прямое следствие для main § 5 (Критерии фальсификации): «нарушение нижней границы Ландауэра» опровергает рамку только в необратимом режиме. Обнаружение системы в Bennett-режиме указывает на границу области определения, не на опровержение программы. Соответствующая операционализация (тест на \(\dot Q > 0\) как признак необратимого режима, перевод в число логических битов при детекции Bennett-asymptotic) --- задача дополнительного исследования, не покрываемая настоящей работой. Bennett-режим как «не служит опровергающим сценарием» зафиксирован в § 5 Блок 1 п. 1 main.



\section{S2.4 Нормировка по когорте: расширенное обсуждение и сигмоидная альтернатива}\label{s2.4-ux43dux43eux440ux43cux438ux440ux43eux432ux43aux430-ux43fux43e-ux43aux43eux433ux43eux440ux442ux435-ux440ux430ux441ux448ux438ux440ux435ux43dux43dux43eux435-ux43eux431ux441ux443ux436ux434ux435ux43dux438ux435-ux438-ux441ux438ux433ux43cux43eux438ux434ux43dux430ux44f-ux430ux43bux44cux442ux435ux440ux43dux430ux442ux438ux432ux430}



\subsection{Расширенное обсуждение референтных когорт}\label{ux440ux430ux441ux448ux438ux440ux435ux43dux43dux43eux435-ux43eux431ux441ux443ux436ux434ux435ux43dux438ux435-ux440ux435ux444ux435ux440ux435ux43dux442ux43dux44bux445-ux43aux43eux433ux43eux440ux442}



Конкретные референтные когорты для шести камертонов раздела 3 фиксируются следующим образом.



\emph{Бактерии} (§ 3.5): когорта прокариотических хемотактических систем, \(n \sim 200\) из BiGG/KEGG-моделей (\emph{E. coli}, \emph{B. subtilis}, \emph{V. cholerae} и подобных).



\emph{Города} (§ 3.6): 120--500 крупных городов из Bettencourt-датасета (Bettencourt et al.~2007).



\emph{Биосфера} (§ 3.7): единственный кейс, для которого нормализация по перцентилю невозможна за отсутствием когорты. Вместо неё используется нормализация по теоретическому верхнему пределу. Биосфера-как-целое нормируется через удельный поток эксергии: NPP \(\sim 10^{14}\) Вт делится на полную живую биомассу \(\sim 5 \cdot 10^{14}\) кг (Bar-On et al.~2018) → удельный \(P_D \sim 0{,}2\) Вт/кг как нормировочный масштаб \(T_v^{\text{ref}}\). Альтернативно: для биосферы-как-целого нормировка через \(N_{\text{max}}\)-эквивалент за окно \(\tau_{\text{year}}\) (см. § 3.7).



\emph{Крупная языковая модель} (§ 3.4): когорта последних поколений foundation-моделей, \(n \sim 50\) от 2020 года.



\emph{Звёзды} (§ 3.1): когорта главной последовательности, \(n \sim 10^4\) из каталога Hipparcos.



\emph{Кристалл и ураган} (§§ 3.2--3.3): нормализация по нулю. Структурный ноль \(T_v = 0\) или \(C_v = 0\) задаётся не положением в когорте, а конструктивным провалом соответствующей ёмкости (отсутствие собственной диссипации, оплачивающей удержание; отсутствие удерживаемой predictive information). Формально: \(\tilde{T}_v(X) := 0\) принимается по конвенции для \(X\) с \(T_v(X) = 0\), и \(\tilde{T}_v(X) := 1\) для \(X\) с \(T_v(X) \ge \max_{\text{cohort}} T_v\); внутри носителя референтной когорты --- формула (2a) основного текста.



\subsection{Сигмоидная альтернатива}\label{ux441ux438ux433ux43cux43eux438ux434ux43dux430ux44f-ux430ux43bux44cux442ux435ux440ux43dux430ux442ux438ux432ux430}



Альтернатива перцентильной нормализации --- сигмоидная нормировка \(\tilde{T}_v = T_v / (T_v + T_v^{\text{ref}})\) через независимый выбор референтного значения \(T_v^{\text{ref}}\) для каждой ёмкости. Сигмоидная процедура даёт абсолютную шкалу и применима к уникальным системам, но требует обоснования \(T_v^{\text{ref}}, C_v^{\text{ref}}, S_v^{\text{ref}}\) из независимых физических соображений. В настоящей работе используется перцентильная процедура. Сигмоидная оставлена как альтернатива для применений, требующих абсолютной шкалы или работы с уникальными системами.



\subsection{\texorpdfstring{Worked example: \(I_v\) для \emph{E. coli} по перцентильной нормализации}{Worked example: I\_v для E. coli по перцентильной нормализации}}\label{worked-example-i_v-ux434ux43bux44f-e.-coli-ux43fux43e-ux43fux435ux440ux446ux435ux43dux442ux438ux43bux44cux43dux43eux439-ux43dux43eux440ux43cux430ux43bux438ux437ux430ux446ux438ux438}



Иллюстративный пошаговый расчёт \(I_v(\textit{E. coli})\) закрывает разрыв между формальным определением (2)--(2a) основного текста и операциональной процедурой. Расчёт даётся как \textbf{разобранный пример, не калиброванное измерение}: числа --- методологическая иллюстрация процедуры, согласованная с заявленным в § 3.5 «\(I_v \sim 0{,}5\) по перцентилю когорты прокариотических хемотактических систем», не самостоятельные эмпирические утверждения.



\emph{(а) Прокси \(T_v\) --- удельный поток эксергии \(P_D \cdot \eta_{\text{ex}}\) в Вт/кг.} Тепловыделение \emph{E. coli} в стационарной фазе \(\sim 1\)--\(3\) пВт/клетку (Liu et al.~2020), сухая масса клетки \(\sim 3 \cdot 10^{-13}\) г, \(P_D \sim 3\)--\(10\) Вт/кг. С эксергетическим КПД \(\eta_{\text{ex}} \sim 0{,}4\) удельный поток эксергии \(P_D \cdot \eta_{\text{ex}} \sim 1\)--\(4\) Вт/кг --- прокси \(T_v(\textit{E. coli})\).



\emph{(б) Распределение прокси по когорте.} Когорта прокариотических хемотактических систем --- \(n \sim 200\) из BiGG/KEGG. Тепловыделение варьирует от \(\sim 0{,}1\) пВт/клетку у автотрофных хемолитотрофов с медленным метаболизмом (\emph{Nitrosomonas}, метаногены) до \(\sim 10\) пВт/клетку у быстрорастущих гетеротрофов (\emph{V. natriegens}). Удельный поток эксергии при сопоставимой массе клетки даёт распределение \(P_D \cdot \eta_{\text{ex}}\) с медианой \(\sim 1\) Вт/кг и межквартильным размахом \([0{,}3;\, 3]\) Вт/кг.



\emph{(в) Ранг }E. coli* в распределении.* При \(P_D \cdot \eta_{\text{ex}}(\textit{E. coli}) \sim 2\) Вт/кг ранг попадает в средне-верхнюю часть распределения, \(\tilde{T}_v \approx 0{,}6\).



\emph{(г) Прокси \(C_v\) и \(S_v\).} \(\tilde{C}_v\) --- нормированная excess entropy \(E\) (Crutchfield and Feldman 2003) или statistical complexity \(C_\mu\) (Shalizi and Crutchfield 2001) на временном ряду метилирования (см. main § 2.3: LZ-сложность принадлежит классу прокси скорости энтропии \(h_\mu\), не \(I_{\text{pred}}\); для \(C_v\) как меры качества модели нужны прокси \(I_{\text{pred}}\)-класса). Для \emph{E. coli} с её адаптивной reactivity к химическим градиентам прокси оценивается несколько выше медианы когорты, \(\tilde{C}_v \approx 0{,}55\) --- иллюстративная точка, формальное распределение excess entropy по когорте \(n \sim 200\) требует отдельной работы. \(\tilde{S}_v\) --- Leiden-модулярность сети хемотактической регуляции (Tar--CheA--CheY--FliM--двигательная аксиально-симметричная сборка); архитектура иерархически-модулярна, \(Q \sim 0{,}5\) по алгоритму Leiden (Traag et al.~2019), что в распределении когорты даёт \(\tilde{S}_v \approx 0{,}4\).



\(\tilde C_v\) и \(\tilde S_v\) для \emph{E. coli} даются как иллюстративные точки в когорте \(n \sim 200\) прокариотических хемотактических систем (BiGG/KEGG); полное распределение excess entropy / Leiden-\(Q\) по этой когорте --- открытая задача operationalisation, не входящая в настоящий разобранный пример. Численная иллюстрация --- методологическая, не калиброванное измерение.



\emph{(д) Геометрическое среднее.} \(I_v = \sqrt[3]{0{,}6 \cdot 0{,}55 \cdot 0{,}4} = \sqrt[3]{0{,}132} \approx 0{,}51\) --- согласуется с заявленным в § 3.5 «\(I_v \sim 0{,}5\)». Чувствительность к выбору прокси внутри одного класса оценивается конвертированием каждого \(\tilde{\cdot}_v\) в диапазон \(\pm 0{,}1\): итоговый \(I_v\) остаётся в коридоре \([0{,}4;\, 0{,}6]\), ширина коридора --- индикатор устойчивости перцентильной процедуры внутри носителя референтной когорты.



Эпистемический статус каждого числа --- иллюстративный. Аналогичный пошаговый расчёт для остальных пяти камертонов требует фиксации когорты с равной дисциплиной, что выходит за рамки настоящей работы; камертонные значения \(I_v\) в § 3 --- порядковые оценки, методологически согласуемые с настоящим разобранным примером.



\section{S3.4 Counterfactual ablation LLM: полный покомпонентный разбор}\label{s3.4-counterfactual-ablation-llm-ux43fux43eux43bux43dux44bux439-ux43fux43eux43aux43eux43cux43fux43eux43dux435ux43dux442ux43dux44bux439-ux440ux430ux437ux431ux43eux440}



Применим counterfactual-процедуру § 2.2 основного текста на конкретном кейсе крупной языковой модели. Кандидаты на включение в границу разбираются последовательно.



\emph{Параметрические веса (статичные после обучения).} Вычитание: модель не отвечает; чистый провал петли. Принадлежат границе генератора, но не принадлежат работающему контуру в момент инференса (веса не обновляются ошибкой).



\emph{KV-cache (память контекста во время выполнения).} Вычитание: модель теряет контекст разговора; за окном \(\tau \sim\) время сессии \(\Delta(-dF_X/dt) > 0\). Принадлежит границе на временах сессии.



\emph{Исполняющий сервер (GPU + ОС).} Вычитание: вычисления невозможны. Принадлежит границе механизма, но энергию подаёт внешняя электросеть.



\emph{Корпорация-владелец (принимающая решения об обучении, развёртывании, оплате).} Вычитание: модель в текущем виде перестаёт обновляться (нет новых версий), но текущая версия продолжает работать на остаточной инфраструктуре. Внешний компонент относительно текущей версии модели; включается в границу при рассмотрении эволюции модельной линейки во времени.



Counterfactual ablation различает четыре нарастающие границы \(B_1 \subset B_2 \subset B_3 \subset B_4\) (см. main § 3.4). При \(B_1\) (только веса) и \(B_2\) (веса + KV-cache + сервер) на окне сессии --- \(\eta_L^{\text{int}}\) \textbf{не определён} через одновременный провал \(T_v^{\text{int}}\) (электросеть внешняя, \(E_{\text{actual}}^{\text{own}} = 0\)) и \(C_v^{\text{int}}\) (нет канала возврата ошибки в работающий контур). При \(B_3\) (обучающий пайплайн: датасет + парк GPU + RLHF-операторы) на окне поколения моделей --- петля принятия решений замыкается на корпоративных операторов, \(\eta_L\) положителен, но субъект --- пайплайн, не модель. При \(B_4\) (корпорация-юр.лицо) --- \(\eta_L^{\text{corp}} \lesssim 10^{-25}\) как \textbf{верхняя граница} через capacity-bound числителя (\(I_{\text{cap}}^{\text{corp}} \lesssim 10^{12}\) бит как верхняя оценка на \(I_{\text{pred}}\), не измеренная predictive information). Counterfactual ablation, таким образом, не даёт LLM витального статуса ни на одной из выбранных шкал. Он операционализирует интуицию, на которой держится диагностика § 3.4 основного текста: расширение границы расширяет ответственного субъекта (модель → пайплайн → корпорация), не растворяет диагноз.



\section{S3.5 Worked example E. coli: подробный анализ чувствительности}\label{s3.5-worked-example-e.-coli-ux43fux43eux434ux440ux43eux431ux43dux44bux439-ux430ux43dux430ux43bux438ux437-ux447ux443ux432ux441ux442ux432ux438ux442ux435ux43bux44cux43dux43eux441ux442ux438}



\subsection{Полные параметры расчёта}\label{ux43fux43eux43bux43dux44bux435-ux43fux430ux440ux430ux43cux435ux442ux440ux44b-ux440ux430ux441ux447ux451ux442ux430}



\emph{Числитель.} Сенсорный канал \(x_t \in \{0, 1\}\) --- бинаризованный знак временного градиента концентрации аспартата (положительный/отрицательный) с шагом \(\delta t \sim 0{,}1\) с (характерное время отклика одного кластера рецепторов). Внутренний канал \(s_t \in \{0, 1, 2, 3, 4\}\) --- число метильных групп на Tar-рецептор (пять дискретных состояний адаптационной методики). Окно прогноза \(\tau \approx 10\) с --- характерное время адаптации системы метилирования. Переходные вероятности \(P(s_{t+\delta t} \mid s_t, x_t)\) задаются аналитически в модели Tu et al.~(2008) (Tu et al.~2008; Shimizu et al.~2010) на стационарном распределении марковского процесса. Альтернативно \(I(s_t; x_{t+\tau})\) оценивается численно через KSG-оценку Kraskov--Stögbauer--Grassberger на симулированных траекториях.



Для одного канала измерения в окне \(\tau\) получается \(I_{\text{pred}} \approx 1\)--\(2\) бит на одно измерение (Mehta and Schwab 2012; Cheong et al.~2011) --- порядок, согласный с информационно-теоретическим анализом биохимического сигналинга. Интегрирование по \(\sim 10^3\)--\(10^4\) измерениям/рецепторам даёт \(10^3\)--\(10^4\) бит за поколение. Типичная \emph{E. coli} несёт \(\sim 7000\) хеморецепторов (Tar+Tsr+Trg+Tap+Aer), организованных в \(\sim 10^3\) кластеров с аллостерическим сопряжением (Monod--Wyman--Changeux-подобная организация). Верхняя граница в предположении независимости каналов --- \(I_{\text{pred}} \sim 10^3\)--\(10^4\) бит за поколение; \textbf{кооперативное сопряжение в кластерах снижает фактическое \(I_{\text{pred}}\) на \(\sim 1\) порядок} (см. main § 3.5), что включено в нижнюю границу интервала чувствительности \([10^{-9}, 10^{-7}]\) итогового \(\eta_L\). Policy-conditioned MI-estimate для хемотактического каскада с учётом кластеров рецепторов --- открытая задача (Tostevin and ten Wolde 2009; Govern and ten Wolde 2014).



\emph{Знаменатель.} Тепловыделение \emph{E. coli} в стационарной фазе роста --- \(\sim 1\)--\(3\) пВт на клетку (Liu et al.~2020). За поколение \(\tau_d \sim 30\) мин \(\approx 1{,}8 \cdot 10^3\) с суммарное тепловыделение \(Q_{\text{heat}} \sim 10^{-9}\) Дж/клетку (порядковая оценка). Точная подстановка \(2 \cdot 10^{-9}\) даёт \(N_{\text{max}} \approx 6 \cdot 10^{11}\), но для согласования с округлёнными числами main § 3.5 здесь принят порядок \(10^{-9}\). По § 2.1 main \(E_{\text{actual}}\) --- подведённая свободная энергия за окно \(\tau_d\), в стационарной фазе равная полной интегральной диссипации. При малой \(\eta_{\text{ATP}}\) на окне поколения калориметрия даёт хорошую аппроксимацию подведённой эксергии: \(E_{\text{actual}} \approx Q_{\text{heat}} \sim 10^{-9}\) Дж/клетку (общий энергобаланс \(E_{\text{actual}} \approx Q_{\text{heat}}/(1-\eta_{\text{ATP}})\) в пределе \(\eta_{\text{ATP}} \to 0\) стационарной фазы). При \(T = 310\) К, \(k_BT\ln 2 \approx 3 \cdot 10^{-21}\) Дж:



\[N_{\text{max}} = \frac{E_{\text{actual}}}{k_B T \ln 2} \sim \frac{10^{-9}}{3 \cdot 10^{-21}} \sim 3 \cdot 10^{11} \;\text{бит на клетку}.\]



Это первая, \emph{эксергическая}, граница: число битов, которые клетка теоретически могла бы стереть, если бы вся её свободная энергия шла на необратимые информационные операции. Реально лишь малая доля эксергии --- по порядку величины \(10^{-3}\), как отношение мощности сенсорно-вычислительной сети (Mehta and Schwab 2012) к полному метаболическому бюджету клетки (Liu et al.~2020) --- направляется на необратимые информационные операции; остальное уходит на синтез биомассы, поддержание ионных градиентов и тепловые потери. Соответствующая \emph{вычислительная} (computational) граница: \(E_{\text{comp}} \sim 10^{-3} \cdot E_{\text{actual}} \sim 10^{-12}\) Дж/клетку и \(N_{\text{max}}^{\text{comp}} = E_{\text{comp}}/(k_B T \ln 2) \sim 3 \cdot 10^{8}\) бит/клетку.



\subsection{Анализ чувствительности}\label{ux430ux43dux430ux43bux438ux437-ux447ux443ux432ux441ux442ux432ux438ux442ux435ux43bux44cux43dux43eux441ux442ux438}



Окно \(\tau\) варьируется в диапазоне \(1\)--\(100\) с --- это сдвигает \(I_{\text{pred}}\) на \(\pm 1\) порядок. На коротких окнах марковская модель даёт меньше информации; на длинных она насыщается из-за ограничения по адаптационной памяти.



Разбиение на интервалы \(x_t\) (бинарное против 4-уровневого) меняет \(I_{\text{pred}}\) на \(\sim 30\%\) --- существенно меньше, чем неопределённость по числу кластеров рецепторов.



Включение всех пяти семейств рецепторов (Tar, Tsr, Trg, Tap, Aer) масштабирует \(I_{\text{pred}}\) примерно линейно с числом сенсоров, оставляя \(\eta_L\) в пределах одного порядка.



Оценка \(\eta_{\text{ATP}}\) (\(< 0{,}1\) на окне поколения по биоэнергетической литературе) меняет \(E_{\text{actual}}\) на \(\sim 10\%\) --- не порядок.



Итоговая консервативная оценка с учётом всех источников неопределённости: \(\eta_L \in [10^{-9}, 10^{-7}]\) через эксергический знаменатель и \(\eta_L^{\text{comp}} \in [10^{-6}, 10^{-4}]\) через вычислительный знаменатель.



\subsection{Узкое место расчёта}\label{ux443ux437ux43aux43eux435-ux43cux435ux441ux442ux43e-ux440ux430ux441ux447ux451ux442ux430}



Узкое место --- оценка \(I_{\text{pred}}\) для полной сенсорной системы: число рецепторов в кластерах варьирует у разных штаммов и условий и берётся здесь из литературных оценок, не из первых принципов. Точная величина требует одновременных MI-измерений на всех известных сенсорных каналах одной популяции, что определяет содержание сопровождающей экспериментальной программы. Полный воспроизводимый расчёт с кодом приведён как часть Supplementary Materials (Python-исходники, сохранённые выводы, документация параметров).



\section{S4.4 Counterfactual ablation термостата: полный покомпонентный разбор}\label{s4.4-counterfactual-ablation-ux442ux435ux440ux43cux43eux441ux442ux430ux442ux430-ux43fux43eux43bux43dux44bux439-ux43fux43eux43aux43eux43cux43fux43eux43dux435ux43dux442ux43dux44bux439-ux440ux430ux437ux431ux43eux440}



Применим counterfactual-процедуру § 2.2 на биметаллическом термостате \(X\) с гистерезисом \(\Delta T_{\text{hyst}} = 1\) К, контролирующем комнату с тепловой инерцией \(\tau_{\text{room}} \sim 30\) мин. Кандидаты на включение в границу \(X\) перебираются последовательно.



\emph{Биметаллическая пластина (датчик и исполнитель).} Вычитание: контур разорван, никакого регулирования. \(\Delta(-dF_X/dt) > 0\) на временах \(< \tau_d\). Принадлежит границе.



\emph{Контактная группа реле (коммутация энергии).} Вычитание: то же --- контур разорван. Принадлежит границе.



\emph{Нагревательный элемент (исполнительное звено, питаемое внешней электросетью).} Вычитание: тепловой источник пропадает; биметалл и реле остаются. Однако \(-dF_X/dt\) оставшейся системы стремится к нулю не через провал петли, а через истощение собственной свободной энергии (биметалл медленно остывает до окружения). Не принадлежит петле принятия решений; внешний компонент.



\emph{Электросеть как источник тока.} Аналогично нагревателю: внешний.



\emph{Уставщик (внешний оператор).} Вычитание: уставка фиксируется, регулирование сохраняется. Внешний.



Граница термостата по counterfactual ablation --- \{биметалл, реле\}. Внутри границы \(\dot{E}_{\text{actual}} \sim 10^{-6}\) Вт --- упругая релаксация металла плюс диссипация реле. За окно \(\tau_d \sim 30\) мин \(\approx 1{,}8 \cdot 10^3\) с: \(E_{\text{actual}} = \dot{E}_{\text{actual}} \cdot \tau_d \sim 2 \cdot 10^{-3}\) Дж. Температура термостата-приёмника --- комнатная, \(T \approx 300\) К; \(k_B T \ln 2 \approx 3 \cdot 10^{-21}\) Дж/бит; \(N_{\text{max}} \sim 6 \cdot 10^{17}\) бит. С \(I_{\text{pred}} \approx 3\) бит для тривиальной задачи «удерживать собственную форму биметалла» при суточном \(\Delta T \sim 10\) К получаем \(\eta_L^{\text{trivial}} \approx 5 \cdot 10^{-18}\). Это центральная оценка; до \(\sim 10^{-17}\) при сезонном \(\Delta T \sim 30\) К и \(I_{\text{pred}} \approx 5\) бит. Внутренний канал, удерживающий predictive information о температуре комнаты, --- деформация биметалла. Оценка: \(I_{\text{pred}} \le \log_2(\Delta T_{\text{room}}/\Delta T_{\text{hyst}}) \approx 3\) бит при суточном \(\Delta T_{\text{room}} \sim 10\) К (до \(\sim 5\) бит при сезонном \(\Delta T_{\text{room}} \sim 30\) К). Ключевое наблюдение: эта \(I_{\text{pred}}\) относится к среде «комната», не к среде «термостат + комната + источник питания». Петля принятия решений, оплачиваемых собственным распадом, у термостата неполна --- нагрев комнаты оплачивается внешней электросетью, не собственной свободной энергией биметалла. Counterfactual ablation даёт \(\eta_L^{\text{thermostat}}\), не определённый в обеих самосогласованных интерпретациях, но через провал разных ёмкостей. В управляющей задаче «регулировать температуру комнаты» проваливается \(T_v^{\text{int}}\): нагрев оплачивает электросеть, \(E_{\text{actual}}^{\text{own}} = 0\) относительно этой задачи. В тривиальной задаче «удерживать собственную форму биметалла» проваливается \(S_v^{\text{int}}\): пассивное thermo-mechanical coupling не образует петли возврата ошибки --- при отклонении формы биметалла никакая собственная диссипация её не корректирует, биметалл просто остаётся в новой форме как функция текущей температуры. Расхождение между интерпретациями указывает не на онтологический переход --- Bruineberg et al.~(2022) формально показывают, что такого перехода без скрытого ввоза эпистемических допущений сделать нельзя. Указывает оно на конвенциональность выбора границы как методологического инструмента: одна и та же физическая система допускает несколько самосогласованных границ, и выбор между ними --- методологическая, не онтологическая дисциплина.



\section{S4.4a Парирование критики Andrews 2021 и Williams 2020}\label{s4.4a-ux43fux430ux440ux438ux440ux43eux432ux430ux43dux438ux435-ux43aux440ux438ux442ux438ux43aux438-andrews-2021-ux438-williams-2020}



Известные критики FEP затрагивают \(\eta_L\) асимметрично.



Andrews (2021) --- «The math is not the territory» --- атакует FEP за невыводимость эмпирического содержания из формализма: без дополнительных гипотез \(F\) можно построить post-hoc для любой системы. В \(\eta_L\) эта критика парируется через привязку числителя и знаменателя к измеримым величинам --- predictive information через прокси \(C_v\), эксергия через термодинамические измерения; свобода post-hoc построения устраняется требованием конкретной операционализации.



Williams (2020) --- «Predictive coding and thought» --- атакует predictive coding как нейронный механизм, не FEP как принцип. Эта критика нейтральна к \(\eta_L\): программа не утверждает predictive coding как механизм, а использует predictive information как формальную меру, методологически независимую от конкретной нейронной архитектуры. Различение predictive coding (Rao and Ballard 1999; Clark 2013 --- нейроархитектурная гипотеза) и predictive information (Bialek et al.~2001 --- теоретико-информационная мера) --- принципиально для последовательной защиты программы от смешанной критики.



\section{S4.7 Organisational closure: autopoiesis, M,R-системы, constraint closure}\label{s4.7-organisational-closure-autopoiesis-mr-ux441ux438ux441ux442ux435ux43cux44b-constraint-closure}



Возражение, заслуживающее отдельного разбора: не является ли self-payment ребрендингом классической линии организационного замыкания --- autopoiesis Матураны--Варелы, M,R-систем Розена, constraint closure Mossio--Montévil--Longo, эпистемического разреза Паттее? Прямой ответ: четыре линии фиксируют разные стороны замыкания на качественном уровне; \(\eta_L\) добавляет к этой традиции одну вещь --- количественную термодинамическую дисциплину через ландауэровский бюджет. Это не альтернатива и не подмена, а отдельный аналитический регистр поверх качественной концептуальной работы.



Паттее (Pattee 1982, 2001) вводит \emph{epistemic cut} как границу между физико-динамическим уровнем и символическим уровнем системы, а \emph{semantic closure} --- как замыкание системы на собственную обработку собственных символов. Биосемиотическая программа Хофмейера (Hoffmeyer 1996) разворачивает эту линию на знаковых процессах живого, фиксируя семиозис как конститутивную черту биологического. Наша конструкция counterfactual ablation границы (§ 2.2 main) --- термодинамическая операционализация этой линии: эпистемический разрез проводится не феноменологически, а интервенционно, через различение компонентов по их вкладу в ландауэровский бюджет; биосемиотическая территория знаковых процессов остаётся за пределами шкалы \(\eta_L\).



Матурана и Варела (Maturana and Varela 1980) определяют автопоэтическую систему как сеть процессов, производящих компоненты, которые рекурсивно реализуют ту же сеть и определяют топологическую границу через собственное производство. Self-payment работает в комплементарном регистре: оплата удержания predictive information физически принадлежит той же системе, которая удерживает информацию. Матурана--Варела описывают organisational closure качественно, через циркулярную структуру «производство компонентов компонентами»; self-payment даёт количественную термодинамическую часть той же дисциплины через ландауэровский бюджет на удержание модели.



Категориальное отличие self-payment от autopoiesis: autopoiesis формализует \textbf{компоненту} замыкания (production of components by components); self-payment формализует \textbf{информационно-термодинамическое тождество} удерживаемого и оплачиваемого. Автопоэтическая система может теоретически удерживать predictive information через внешний термодинамический подвод (поддержание клеточной культуры in vitro: автопоэтична метаболически, но удерживается внешним подводом); \(\eta_L\) снимет ей витальный статус через знаменатель. Это конкретный категориальный сдвиг, не только количественный.



M,R-системы Розена (Rosen 1991) --- категорно-теоретическая абстракция замыкания относительно efficient causation: организм есть система, в которой каждая причинная роль (Aristotle's efficient cause) замкнута внутри сети. Self-payment \(\eta_L\) --- её термодинамическое прочтение: замыкание реализуется через физическую петлю модель↔диссипация, а не через категорную абстракцию. Важное содержательное различение: Розен в \emph{Life Itself} аргументирует, что closure-to-efficient-causation не-вычислимо и не-формализуемо в машинно-теоретических терминах. \(\eta_L\) не противоречит этому утверждению. Он не утверждает, что биология сводима к вычислению; он утверждает, что термодинамическая цена удержания модели --- необходимое условие витальности, измеримое формально и не претендующее на исчерпание феномена жизни. Шкала фиксирует один обязательный термодинамический срез, оставляя открытым вопрос о невычислимых аспектах организационного замыкания.



Constraint closure Mossio--Montévil--Longo (Mossio et al.~2016) фиксирует closure of biological constraints: каждое биологическое ограничение замыкается через производство другого ограничения, и сеть таких ограничений замыкается в себе. Self-payment работает на уровне информационно-энергетической петли (predictive information ↔ ландауэровская диссипация), а не на уровне сети ограничений; конструкции ортогональны, не альтернативны. Constraint closure описывает топологию взаимного производства ограничений; self-payment даёт термодинамическую плату, которую эта топология должна оплачивать своим собственным распадом.



\emph{Дифференциальный тест шкалы относительно organisational closure.} Рассмотрим пару WT \emph{E. coli} / минимальная клетка JCVI-syn3.0. Обе системы автопоэтичны по Матуране--Вареле, обладают M,R-замкнутостью по Розену и constraint closure по Mossio--Montévil--Longo: качественные критерии organisational closure не различают их. Шкала \(\eta_L\) даёт дифференциальный результат: \(\eta_L^{\textit{E. coli}}\) и \(\eta_L^{\text{syn3.0}}\) различаются на 1--2 порядка через \textbf{оба компонента}. Знаменатель: \(E_{\text{actual}}^{\text{syn3.0}}\) существенно меньше (меньший метаболический поток \textasciitilde473 гена против \textasciitilde4300). Числитель: \(I_{\text{pred}}^{\text{syn3.0}}\) меньше (резко урезанный сенсорный репертуар, отсутствие двухкомпонентных хемотактических каскадов). \(\eta_L\) ранжирует системы внутри класса organisational closure через произведение информационной и термодинамической компонент; конкретный знак отношения \(\eta_L^{\textit{E. coli}} / \eta_L^{\text{syn3.0}}\) --- иллюстративная оценка по доступной литературе (Hutchison et al.~2016, для генома; метаболическая калориметрия для \(E_{\text{actual}}\)), требующая независимой эмпирической верификации.



Свод по линии замыкания. Четыре концептуальных предшественника фиксируют различные аспекты замыкания: epistemic-symbolic у Паттее, production-component у Матураны--Варелы, efficient-causation у Розена, constraint-network у Mossio--Montévil--Longo. \(\eta_L\) добавляет к этой традиции количественную термодинамическую дисциплину: ландауэровский бюджет на удержание predictive information в собственных степенях свободы системы. Это не ребрендинг четырёх линий, а построение количественной шкалы поверх их качественной концептуальной работы. Позиционное соотношение симметрично шести линиям § 4.6 main: каждая классическая программа операционализирует одну сторону феномена; \(\eta_L\) собирает термодинамическую плату в одну измеримую величину, не претендуя на замену соответствующей качественной аналитики.



\section{S4.8 Линия происхождения жизни и концептуальные предшественники}\label{s4.8-ux43bux438ux43dux438ux44f-ux43fux440ux43eux438ux441ux445ux43eux436ux434ux435ux43dux438ux44f-ux436ux438ux437ux43dux438-ux438-ux43aux43eux43dux446ux435ux43fux442ux443ux430ux43bux44cux43dux44bux435-ux43fux440ux435ux434ux448ux435ux441ux442ux432ux435ux43dux43dux438ux43aux438}



Линия организационного замыкания (§ S4.7) фиксирует синхронный аспект витальности --- как уже сформированная живая система удерживает себя. Параллельно ей развивается современная программа исследования \textbf{происхождения жизни} (origin of life, OoL), фиксирующая диахронический аспект --- при каких условиях замкнутая петля впервые возникает в геохимической системе. В Theory in Biosciences и смежных журналах эта программа представлена пятью относительно автономными традициями; каждая операционализирует одну сторону перехода абиотическое→биотическое, и каждая является концептуальным предшественником одного из срезов \(\eta_L\)-постулата. Подраздел артикулирует это соотношение явно --- не как заявку на закрытие проблемы OoL, а как разметку места \(\eta_L\) внутри уже существующего поля.



\emph{Гиперцикл Эйгена--Шустера.} Гиперцикл (Eigen 1971; Eigen and Schuster 1979) --- автокаталитическое замыкание уровня репликаторов с информационным шаблон-носителем. Множество репликаторов организовано в каталитический цикл, в котором каждый член катализирует репликацию следующего; замыкание петли возврата ошибки происходит через перекрёстно-каталитические зависимости между членами цикла. \textbf{Error catastrophe} (Eigen 1971) фиксирует порог между сохранением и потерей информации в репликаторной петле при росте скорости мутаций. Этот порог концептуально гомологичен структурному переходу \(\eta_L = 0 \to \eta_L > 0\) через постулат стационарного учёта (§ 2.1 main, шаг iv): и там, и там речь о пороге, ниже которого информационно-несущая структура теряет способность удерживать собственную модель против шумового канала. Гиперцикл даёт доклеточный разворот числителя \(\eta_L\): \(I_{\text{pred}}\) кодируется в репертуаре репликаторов, \(T_v\) --- в каталитическом обороте гиперцикла, замыкание петли возврата ошибки --- через перекрёстно-каталитические зависимости. \(\eta_L\) дополняет гиперцикл количественной ландауэровской дисциплиной: к структурному условию репликаторного замыкания с шаблон-носителем добавляется требование, чтобы свободная энергия, удерживающая эту замкнутость, физически оплачивалась собственным каталитическим оборотом цикла, а не подвергалась чисто внешней накачке.



\emph{RAF-сети Хордейка--Стила.} Формализм reflexively autocatalytic and food-generated (RAF) сетей (Hordijk and Steel 2004; Hordijk 2019) описывает условия, при которых множество химических реакций образует коллективно автокаталитическую замкнутую сеть над заданным пищевым множеством. Каждая реакция сети катализируется хотя бы одним её собственным продуктом. Вся сеть может быть собрана из пищевого множества за конечное число шагов. RAF --- структурный предшественник композита \(C_v\) (информационная ёмкость через топологию каталитической сети) и \(S_v\) (системная ёмкость как иерархическая модулярность): коллективная автокаталитичность есть структурное условие появления замкнутой петли. Однако RAF --- критерий \textbf{замкнутости сети}, не критерий, \textbf{кто оплачивает её удержание} в смысле § 2.1 main. RAF-сеть может удовлетворять условиям замкнутости и в то же время оплачиваться внешним химическим потоком, не оплачиваться собственной диссипативной работой системы. \(\eta_L\) дополняет RAF термодинамической дисциплиной self-payment: к структурному условию автокаталитической замкнутости добавляется требование, чтобы свободная энергия, удерживающая эту замкнутость, физически принадлежала самой сети.



\emph{Dissipative adaptation Энглэнда.} Термин «dissipative adaptation» введён в Perunov, Marsland, and England (2016) «\emph{Statistical physics of adaptation}» (Phys. Rev.~X 6:021036); его статистико-физическая основа --- England (2013) «\emph{Statistical physics of self-replication}» (J. Chem. Phys. 139:121923), в которой выведена нижняя граница на тепловыделение для самореплицирующихся систем (через темп роста, внутреннюю энтропию, долговечность). Формализм Перунова--Марсланда--Энглэнда описывает dissipative adaptation как процесс, в котором в системах вдали от равновесия, подвергающихся внешней движущей силе, статистически фиксируются макросостояния, которые особенно эффективно поглощают и диссипируют энергию из среды. Это прямой термодинамический предшественник \textbf{знаменателя} \(\eta_L\) --- \(T_v\)-стороны программы (термодинамическая ёмкость как способность системы участвовать в неравновесном обмене). Однако dissipative adaptation фиксирует адаптацию диссипативной структуры как таковой, не предсказание о среде, удерживаемое этой структурой: в формализме Энглэнда нет числителя \(I_{\text{pred}}\). \(\eta_L\) добавляет к диссипативной адаптации требование, чтобы адаптированное состояние (adapted state) удерживало собственную модель среды --- переходя от чистой диссипативной структуры к структуре, оплачивающей собственное моделирование.



\emph{Set-inclusion \{P-M-E adapted\} ⊃ \{\(\eta_L\) vital\}.} Программа Перунова--Марсланда--Энглэнда охватывает шире, чем \(\eta_L\)-витальность: любая система далеко от равновесия, надёжно поглощающая и диссипирующая энергию внешнего driving field, квалифицируется как «physical adaptation» --- независимо от наличия петли самооплаты. Перунов и соавторы эксплицитно: «there may be many examples of `well-adapted' structures that did not have parents» (P-M-E 2016, стр. 21) --- physical adaptation как universal phenomenon, отдельное от Darwinian. Эмпирическая иллюстрация самой P-M-E 2016 --- серебряные наностержни в световом поле (Ito et al.~2013), которые статистически «выучивают» согласование поверхностных плазмонных резонансов с частотой накачивающего поля; примеры подобной адаптированной, но нерепликаторной динамики также фиксируются в полярных актиновых паттернах и активных жидких кристаллах. В \(\eta_L\)-терминах эти системы --- \textbf{adapted, но не vital}: они проходят P-M-E «physical adaptation» test (надёжное согласование поглощения-диссипации с внешней накачкой), но проваливают counterfactual ablation § 2.2 --- оплата извне (световое поле, приток ATP), нет петли self-payment. Аналогично термостат § 4.4: P-M-E-адаптирован к питанию от электросети, но \(\eta_L^{\text{int}}\) не определён через одновременный провал \(T_v^{\text{int}}\) и \(C_v^{\text{int}}\). Set-inclusion формализует разделение: P-M-E adaptation --- \textbf{необходимое, но не достаточное} условие \(\eta_L\)-витальности. \(\eta_L\) выделяет внутри P-M-E-adapted класса subset с термодинамической замкнутостью бюджета self-payment.



\emph{Major transitions in evolution.} Программа Maynard Smith and Szathmáry (1995) --- обновлено в Szathmáry (2015) с явным включением информационно-эволюционных аргументов --- фиксирует последовательность ключевых эволюционных переходов. Это переходы от молекулярных репликаторов к хромосомам и клеткам, от прокариот к эукариотам, от одноклеточных к колониальным и многоклеточным, от индивидов к социальным группам. Каждый из них вводит \textbf{новый уровень субъекта эволюционного процесса} с переадресацией репликации и ошибочной нагрузки. В \(\eta_L\)-терминах major transition есть событие появления \textbf{нового субъекта self-payment}: после перехода свободная энергия, удерживающая predictive information о среде, замыкается на нового агрегатного актора (клетка как актор поверх молекул, многоклеточный организм как актор поверх клеток). Соответствие открывает исследовательскую программу: квалифицируется ли каждое major transition по Maynard Smith--Szathmáry как фазовый переход в субъекте \(\eta_L\) --- с собственным числителем \(I_{\text{pred}}^{\text{new}}\), оплачиваемым диссипацией на новом уровне организации? Этот вопрос предъявляет конкретный термодинамический критерий на иерархию эволюционных уровней и остаётся открытой задачей будущей работы.



\emph{Геохимические программы OoL.} Программы Pross (2012), Smith and Morowitz (2016) и Кауфмана (Kauffman 2019) рассматривают переход abiotic→biotic как обнаружение замкнутой петли поверх геохимического потока. Гидротермальные системы, щелочные источники, метаболические катализаторы образуют контур, в котором часть свободной энергии внешней геохимической движущей силы начинает удерживаться внутри собственной структуры контура. Pross формулирует это как переход от термодинамической стабильности к динамической кинетической устойчивости. Smith--Morowitz --- как фазовый переход в геохимической сети с появлением метаболического ядра (metabolic core). Кауфман --- как обнаружение соседнего возможного (adjacent possible) через спонтанное накопление каталитических замыканий. Alkaline-vent OoL программа Martin--Russell (Martin and Russell 2003) фиксирует геохимический переход к хемоавтотрофным прокариотам как параллельную линию относительно метаболизм-first Smith--Morowitz и Pross. Walker (2017b) в обзоре \emph{Reports on Progress in Physics} картирует OoL как задачу физики, замыкая контекст этих программ как ландшафта.



В \(\eta_L\)-терминах все три геохимические программы (Pross, Smith--Morowitz, Kauffman) артикулируют один и тот же фазовый переход. Момент origin of life --- момент, когда геохимическая система переходит от \emph{внешней оплаты} (полная свободная энергия проходит сквозь систему транзитом, не удерживаясь) к \emph{self-payment}. Часть свободной энергии замыкается внутри собственного контура и физически оплачивает удержание модели среды. Эта формулировка имеет статус риск-предсказания для будущего расширения методологии за пределы стационарного режима; в рамках настоящей работы операциональная процедура измерения \(\eta_L\) на pre-biotic геохимическом субстрате не определена и составляет открытую задачу. Counterfactual ablation границы § 2.2 main и MI-оценка \(I_{\text{pred}}\) через сенсорный канал применимы к уже-сформированным биологическим системам. Их перенос в pre-biotic регистр требует разработки proxies для \(I_{\text{pred}}\) и self-payment в эволюционирующей геохимической сети.



\emph{Параллельная программа Chirumbolo--Vella 2026.} Chirumbolo and Vella (2026) в \emph{Theory in Biosciences} формулируют origin of life как информационно-диссипативный процесс с water-mediated informational entropy gradients. Их рамка --- концептуальный сосед нашей через общую онтологию informational dissipation; различие --- в операциональной мере: они оставляют шкалу качественной, \(\eta_L\) задаёт количественную ландауэровскую дисциплину поверх той же интуиции о self-organizing dissipative systems.



\emph{Сводный абзац.} Пять программ --- гиперцикл Эйгена--Шустера, RAF-сети, dissipative adaptation, major transitions, геохимические OoL --- каждая операционализирует одну сторону феномена возникновения и поддержания живых систем. Соответствующие стороны: replicator-level автокаталитическое замыкание с template-носителем (hypercycle), структурная замкнутость каталитической сети (RAF), диссипативная адаптация, эволюционный субъект, геохимический фазовый переход. \(\eta_L\) не конкурирует ни с одной из них и не претендует на закрытие проблемы происхождения жизни; вместо этого шкала объединяет их в единую безразмерную меру через постулат self-payment. Это термодинамическая операционализация условия, которое каждая из пяти программ артикулирует со своей стороны: положительность \(\eta_L\) --- необходимое условие момента origin of life в любой из них. Дальнейшее уточнение каждого из соответствий --- открытая исследовательская задача.



\section{S4.9 Позиционирование в современных демаркационных дебатах}\label{s4.9-ux43fux43eux437ux438ux446ux438ux43eux43dux438ux440ux43eux432ux430ux43dux438ux435-ux432-ux441ux43eux432ux440ux435ux43cux435ux43dux43dux44bux445-ux434ux435ux43cux430ux440ux43aux430ux446ux438ux43eux43dux43dux44bux445-ux434ux435ux431ux430ux442ux430ux445}



Помимо линий организационного замыкания (§ S4.7) и происхождения жизни (§ S4.8), \(\eta_L\)-программа явно соотносится с четырьмя устойчивыми линиями современной philosophy of biology, фиксирующими демаркационную трудность концепта «жизнь». Эти линии артикулированы в § 1 main первым параграфом; здесь --- позиционная диагностика, как каждая из них соотносится с операциональной шкалой \(\eta_L\). Главный методологический тезис подраздела: \(\eta_L\) не предлагает очередное определение жизни, а вносит в продолжающиеся дебаты формальную количественную шкалу, на которой обозначенные линии могут совместно работать.



\emph{Anti-definitionism Cleland.} Cleland (2019) аргументирует, что поиск определения жизни через списки признаков методологически бесплоден и десятилетиями не приводит к сходимости. Вместо словарного определения нужна «theory of life» --- теоретическая рамка, в которой жизнь занимает положение, определяемое объяснительной структурой, а не списочной демаркацией. Эмпирический фон линии --- реестр Trifonov (2011), зафиксировавший лексический анализ 123 определений жизни без сходимости к общему ядру; anti-definitionism Cleland --- методологическая реакция на этот материал. \(\eta_L\)-программа \emph{присоединяется} к этому диагнозу: определения через признаки ломаются при пересечении границ --- биосфера, LLM-corp, город, минимальная клетка предъявляют разные комбинации признаков, не подводимые под единый ассортативный реестр. \(\eta_L\) --- не очередное определение, а операциональный дискриминатор: безразмерная шкала, на которой каждая система получает положение через измеримое отношение predictive information к ландауэровскому бюджету, без претензии исчерпать феномен жизни. Различие. Anti-definitionism Cleland декларирует методологическую бесплодность критерия-через-признаки; \(\eta_L\) предлагает не определение через признаки, а операциональный дискриминатор с явным режимом фальсификации. Это не theory of life в полном смысле, а одна из обязательных размерностей такой теории.



\emph{Cluster definitions Mariscal--Doolittle.} Программа «life and life only» Mariscal and Doolittle (2020) рассматривает жизнь как кластер связанных концепций, не сводимых к единственному определению. Разные подходы освещают разные стороны феномена; попытка свести их к одной формуле теряет аналитическую содержательность. Эта позиция содержательно совместима с \(\eta_L\) как process-property. Положительность \(\eta_L > 0\) не есть сущностное определение «жизни как таковой», а свойство процесса самооплаты predictive information, реализуемое \textbf{множественно} --- химической петлёй бактерии, биогеохимическим контуром биосферы, синтетическим minimal cell, гипотетической небиологической архитектурой с замкнутой петлёй модель↔диссипация. Кластерная природа концепта сохраняется; \(\eta_L\) вносит в кластер одну общую термодинамическую размерность. Различие: cluster-definitions допускают, что не все жизнеподобные феномены имеют единый термодинамический инвариант; \(\eta_L\) предлагает один единый термодинамический инвариант для синхронной размерности витальности. Программа не претендует на унификацию диахронных (популяционно-эволюционных) аспектов (§ 2.7 main). Это рискованное утверждение, фальсифицируемое наблюдением жизнеподобной системы по трём независимым синхронным каналам с устойчивым \(\eta_L = 0\) (§ 5 main).



\emph{Operational definitions at frontiers Bich--Green.} Bich and Green (2018) прямо обсуждают, как операциональные критерии должны вести себя на границах биологии --- на вирусах, протоклетках, синтетических конструкциях. \(\eta_L\) --- операциональный критерий ровно этой формы. Для вирусной системы, не удерживающей замкнутой петли self-payment и паразитирующей на хозяине, числитель и знаменатель определены так, что собственная predictive information вируса не оплачивается собственной диссипацией хозяина-независимо. Шкала даёт \(\eta_L = 0\) для вируса-в-изоляции и согласована с биологической конвенцией, которая обычно не присваивает вирусу автономную витальность. Этот пункт --- методологический, а не феноменологический: \(\eta_L\) ведёт себя в пограничных случаях так, как требует от операционального определения программа Bich--Green, не воспроизводя дилеммы определений через списки. Различие: на вирусной системе \(\eta_L = 0\) через числитель; вирус-в-хозяине меняет статус --- программа должна сказать, \textbf{чей} \(\eta_L\) растёт при заражении (переадресация объекта, § 2.2 main). Bich--Green оставляют вопрос конвенциональности границы вируса открытым; \(\eta_L\) фиксирует через counterfactual ablation.



\emph{Process biology Pradeu.} Pradeu (2018) аргументирует, что организм методологически продуктивнее мыслить как \textbf{процесс}, не как объект --- биологическая индивидуальность задаётся через временную динамику, не через статическую объектную идентичность. \(\eta_L\) операционализирует ровно это: шкала вычисляется на окне \(\tau_d\) устойчивости и удерживается за окно (§ 2.6 main); витальность определена не как свойство системы-как-объекта, а как характеристика динамики системы за временное окно. Окно \(\tau_d\), counterfactual ablation границы и стационарность триады инвариантов (§ 2.3 main) --- все три конструкции операционализируют процессуальную онтологию Pradeu в термодинамическом регистре. Различие: процессный взгляд Pradeu исключает идентифицируемый субстрат во все моменты времени; \(\eta_L\) требует фиксации границы \(X\) через counterfactual ablation на каждом окне \(\tau_d\). Это не введение объектности обратно --- counterfactual ablation не предполагает постоянство объекта, а только локальную причинно-избирательную замкнутость на окне.



\emph{Информационно-теоретическая индивидуальность Krakauer et al.} Информационно-теоретическая операционализация индивидуальности Krakauer et al.~в самом \emph{Theory in Biosciences} (Krakauer et al.~2020) предлагает разложение системы-среды через каналы пропаганды информации с тремя формами индивидуальности (organismal/colonial/driven). \(\eta_L\) комплементарно их рамке: их аппарат фиксирует информационную сторону демаркации без термодинамической платы; self-payment добавляет ландауэровский знаменатель, делая трёхформную типологию зависимой от того, \textbf{кто оплачивает} удержание информации.



\emph{Basic autonomy Ruiz-Mirazo--Moreno.} Программа basic autonomy Ruiz-Mirazo and Moreno (2004) --- линия, соединяющая autopoiesis с термодинамической constraint-логикой ближе всех к self-payment. Различие: \(\eta_L\) позиционируется как термодинамическая дисциплинация поверх basic autonomy через явное тождество границ модели и оплачивающей её диссипации.



\emph{Сводный абзац.} \(\eta_L\)-программа методологически совместима с anti-definitionist линией Cleland, cluster-definition программой Mariscal--Doolittle, operational-frontiers подходом Bich--Green, process-biology позицией Pradeu, информационно-теоретической рамкой Krakauer et al.~и basic-autonomy программой Ruiz-Mirazo--Moreno --- не как альтернатива им, а как формальная количественная шкала, на которой эти линии могут совместно работать. Это вклад в продолжающиеся демаркационные дебаты, не предложение нового определения «что есть жизнь». Программа --- операциональный дискриминатор, удерживающий процессуальную онтологию, реализуемый множественно через разные прокси числителя и знаменателя, и ведущий себя на границах в согласии с биологической конвенцией.



\section{S5 Критерии фальсификации: полные операциональные протоколы}\label{s5-ux43aux440ux438ux442ux435ux440ux438ux438-ux444ux430ux43bux44cux441ux438ux444ux438ux43aux430ux446ux438ux438-ux43fux43eux43bux43dux44bux435-ux43eux43fux435ux440ux430ux446ux438ux43eux43dux430ux43bux44cux43dux44bux435-ux43fux440ux43eux442ux43eux43aux43eux43bux44b}



Main § 5 содержит сжатые формулировки восьми линий фальсификации, объединённых в три блока. Здесь --- полные операциональные протоколы с численными критериями, p-values, размерами выборок, требованиями предрегистрации и условными оговорками. Структура S5.1--S5.3 параллельна Блокам 1--3 main § 5; S5.4 содержит сводный лакатошевский паспорт программы.



\emph{Примечание о нумерации.} Девять пронумерованных пунктов трёх блоков редуцируются к восьми содержательно независимым линиям. Пункт 1 Блока 2 (статистический аспект --- нарушение монотонности) и пункт 1 Блока 3 (конструктивный аспект --- степенная зависимость, Прогноз 1) --- две формулировки одной риск-ставки о темпе эволюционной адаптации.



\subsection{S5.1 Блок 1: прямые опровержения шкалы (операциональные детали)}\label{s5.1-ux431ux43bux43eux43a-1-ux43fux440ux44fux43cux44bux435-ux43eux43fux440ux43eux432ux435ux440ux436ux435ux43dux438ux44f-ux448ux43aux430ux43bux44b-ux43eux43fux435ux440ux430ux446ux438ux43eux43dux430ux43bux44cux43dux44bux435-ux434ux435ux442ux430ux43bux438}



\emph{Блок 1 п. 1 --- \(T_v\) провал (полный протокол).} Наблюдение: стационарное удержание \(I_{\text{pred}} > 0\) при \(E_{\text{actual}}\) ниже ландауэровской границы \(I_{\text{pred}} \cdot k_B T \ln 2\) в необратимом режиме реализации. Операциональный тест необратимости: измеримое тепловыделение \(\dot Q > 0\) при выполнении информационной операции (обратимая Bennett-реализация даёт \(\dot Q \to 0\) в пределе медленного протокола). Фиксация \(T\): температура термостата, в который система фактически сбрасывает \(E_{\text{actual}}\) --- не среды-источника и не уставки управления. Такое наблюдение реализует \(\eta_L > 1\); опровергает либо постулат стационарного учёта (в) Утверждения § 2.1 main, либо одно из условий (а)--(б); устраняет \(\eta_L \le 1\) как границу шкалы и обнуляет \(T_v\) как необходимое условие.



\emph{Bennett-режим: оговорка.} Bennett-режим обратимых вычислений (Bennett 1973, 1982) не служит опровергающим сценарием: он вне области определения и требует отдельной операционализации через число логических битов. Bennett-режим --- предсуществующее физическое ограничение принципа Ландауэра, осознанное задолго до настоящей программы (Bennett 1973); его исключение не специфично для \(\eta_L\)-программы (любая ландауэровская мера сталкивается с тем же вопросом области определения). Опровержение через Bennett-демонстрацию засчитывается, если автор контр-примера сначала операционализирует «число логических битов» как канонический числитель и покажет \(\eta_L^{\text{log}} > 1\).



\emph{Блок 1 п. 2 --- \(C_v\) провал (полный протокол).} Наблюдение: система с положительным \(C_v^{\text{op}}\) при \(\dot{H}_{\text{erasure}} = 0\) на окне \(\tau_d\) --- пример системы, удерживающей predictive information без необходимости платы ландауэровского стирания за каждый predictive бит. Обнуляет \(C_v\) как необходимое условие.



\emph{Блок 1 п. 3 --- \(S_v\) провал (полный протокол).} Наблюдение: система с замкнутой петлёй self-payment (\(\eta_L > 0\) по counterfactual ablation § 2.2 main) при одновременном выполнении двух условий: \(Q < Q_{\text{null}} + 2\sigma\) (модулярность не отличима от degree-preserving нуль-модели --- configuration model, сохраняющего распределение степеней) и \(C(k) \not\sim k^{-\beta}\) с \(\beta \in [0{,}5; 1{,}5]\) (отсутствие иерархической сигнатуры по Ravasz--Barabási 2003). Требуется ансамбль \(\ge 10\) независимых биологических субъектов. Опровержение засчитывается при \(p < 0{,}01\) против null-гипотезы случайной топологии при заявленной мощности и предрегистрации протокола. Обнуляет \(S_v\) как необходимое условие.



\subsection{S5.2 Блок 2: тесты прокси и операционализации (операциональные детали)}\label{s5.2-ux431ux43bux43eux43a-2-ux442ux435ux441ux442ux44b-ux43fux440ux43eux43aux441ux438-ux438-ux43eux43fux435ux440ux430ux446ux438ux43eux43dux430ux43bux438ux437ux430ux446ux438ux438-ux43eux43fux435ux440ux430ux446ux438ux43eux43dux430ux43bux44cux43dux44bux435-ux434ux435ux442ux430ux43bux438}



\emph{Блок 2 п. 1 --- нарушение монотонности (полный протокол).} Наблюдение: систематическое нарушение монотонной зависимости темпа эволюционного обучения от \(\eta_L\). Если в контролируемом эксперименте по экспериментальной эволюции (бактериальные культуры в хемостате при варьировании температуры, плотности популяции и источника углерода) системы с \(\eta_L\), отличающимся на порядок, дают темпы адаптации, неразличимые статистически на достаточно длинных окнах, центральное отношение лишается числового смысла. Конструктивная проверка совпадает с протоколом S5.3 (Прогноз 1).



\emph{Блок 2 п. 2 --- витальность по трём каналам при \(\eta_L = 0\) (полный протокол).} Наблюдение: система, указывающая на витальность по трём независимым каналам --- темп адаптации, скорость самовосстановления, независимые биохимические маркеры самопроизводства (в смысле дарвиновской эволюции; ср. NASA-определение жизни как «химической системы, способной к дарвиновской эволюции» (Joyce 1994) --- здесь оно привлекается как конвенциональный внешний референт, не как конкурирующая теория витальности; смена субстрат-связанного регистра на этом шаге намеренная). При этом \(\eta_L = 0\) при всех операционализациях \(I_{\text{pred}}\) и \(E_{\text{actual}}\) через стандартные домен-специфичные прокси. Это независимый от \(\eta_L\) критерий: темп адаптации, скорость самовосстановления и биохимические маркеры самопроизводства операционализуются через измеримые величины, не сводимые к числителю \(\eta_L\). Систематическое расхождение указало бы либо на дефект числителя как операционального прокси для предсказательной информации (методологический уровень фальсификации, см. § 2.7 main), либо на несостоятельность самого концепта витальности через эффективность самомоделирования (онтологический уровень).



\emph{Блок 2 п. 3 --- конвенционально живая система с \(\eta_L = 0\) (полный протокол).} Наблюдение: биологическая система, заведомо живая по конвенциональным критериям, у которой \(\eta_L\) строго равен нулю при всех разумных операционализациях. Опровергает адекватность шкалы как операционального дискриминатора витальности.



\subsection{S5.3 Блок 3: прогрессивные предсказания (операциональные детали)}\label{s5.3-ux431ux43bux43eux43a-3-ux43fux440ux43eux433ux440ux435ux441ux441ux438ux432ux43dux44bux435-ux43fux440ux435ux434ux441ux43aux430ux437ux430ux43dux438ux44f-ux43eux43fux435ux440ux430ux446ux438ux43eux43dux430ux43bux44cux43dux44bux435-ux434ux435ux442ux430ux43bux438}



\emph{Блок 3 п. 1 --- Прогноз 1, степенная зависимость темпа адаптации (полный протокол).} Темп убыли расхождения внутренней модели со средой, нормированный на сложность среды, должен скейлироваться с \(\eta_L\) как монотонно растущая функция. Конкретный протокол: культура \emph{E. coli} в хемостате с управляемым потоком углерода (глюкоза, лактоза, ацетат); \(\eta_L\) варьируется через температуру (25--42 °C, пять уровней) и плотность популяции (три уровня); адаптация --- убыль расхождения между кодируемым уровнем метилирования и наблюдаемым градиентом за 50 поколений. Положительное предсказание: коэффициент корреляции (Пирсон между \(\log r_{\text{adapt}}\) и \(\log \eta_L\)) \(\rho \in [0{,}4;\, 0{,}7]\) на 45 линиях; \(r_{\text{adapt}} = A \cdot \eta_L^{\alpha}\) с \(\alpha \in [0{,}3;\, 1{,}0]\). Диапазон узок относительно нулевой гипотезы \(\alpha = 0\) и относительно гипотезы \(\alpha \ge 2\) (сильная нелинейность); ожидаемая независимая точка проверки --- LTEE-данные Ленски, не использованные для выведения нижней границы. Верхняя граница \(\alpha = 1\) --- прямая пропорциональность темпа адаптации потоку \(I_{\text{pred}}\); нижняя \(\alpha = 1/3\) выводится из структуры (2) main: \(I_v\) --- кубический корень произведения трёх ёмкостей, минимальная зависимость \(r_{\text{adapt}} \propto \sqrt[3]{\eta_L}\) соответствует вкладу одной ёмкости из трёх. Шкала калибрована на разобранном примере § 3.5 main (\(\eta_L^{\text{comp}} \in [10^{-6}, 10^{-4}]\)). Опровержение: \(\hat{\alpha}\) значимо вне \([0{,}3;\, 1{,}0]\) при \(p < 0{,}01\) (неправильный показатель степени); \(\hat{\rho} < 0{,}2\) при заявленной мощности (нет detectable trend); либо \(\hat{\rho} < 0\) (тренд в противоположном направлении --- качественное опровержение).



\emph{Блок 3 п. 2 --- Прогноз 2, монотонность по окну согласования (полный протокол).} В режиме приближения к стационарности, когда \(I_{\text{pred}}\) ещё не насыщен и накапливается быстрее интегральной диссипации, стационарная формулировка предсказывает: при нормировке на единичную скорость диссипации (\(\dot{E}_{\text{actual}} = \text{const}\) за пред-стационарный транзиент с окном \(\Delta\tau \ll \tau_d\)) ожидается монотонный рост \(\eta_L(\Delta\tau)\) за интервал согласования. Проверка на тех же 45 линиях \emph{E. coli}: измерять \(\eta_L(\Delta\tau)\) для \(\Delta\tau\) от часа до поколения с независимой фиксацией \(\dot{E}_{\text{actual}}\). Убывание \(\eta_L\) с \(\Delta\tau\) в этом режиме --- контр-аргумент к стационарной рамке как таковой, не делегирование в сопровождающую работу.



\emph{Блок 3 п. 3 --- Прогноз 3, структурный порог \(f^*\) (полный протокол).} Рассматривается класс гипотетических автономных воплощённых агентов с фотовольтаическим или аккумуляторным энергоснабжением и физическим актуатором с собственным термодинамическим бюджетом. Для этого класса предсказывается структурный порог \(f^*\) доли интернализации энергобюджета (отношение \(E_{\text{actual}}^{\text{own}}/E_{\text{actual}}^{\text{total}}\)). Ниже порога counterfactual ablation § 2.2 main даёт \(\eta_L^{\text{int}}\) не определён (структурный провал петли самооплаты); выше --- измеримый \(\eta_L^{\text{int}} > 0\). Числовая локализация \(f^* \in [0{,}1;\, 0{,}5]\) --- исследовательская гипотеза программы; формальное derivation границ через counterfactual ablation --- открытая задача. По двухуровневой формулировке § 3.4 main структурный предикат counterfactual ablation \textbf{бинарен} (самооплачиваемый канал возврата ошибки есть либо нет), тогда как \(f^*\) --- \textbf{эмпирический} порог на непрерывном прокси \(f\), выше которого замыкание петли статистически устойчиво; \(f^*\) --- регулярность совместной интернализации, не сама структурная решающая граница.



\emph{Оговорка об области применимости.} Предсказание неприменимо к текущим cloud-hosted LLM без сохранения состояния, где \(E_{\text{actual}}^{\text{own}} = 0\) конструктивно (электросеть и инфраструктура внешние относительно прямого прохода), \(f = 0\) по построению.



\emph{Дифференциальное предсказание относительно AI-safety литературы.} AI-safety программы (Hubinger et al.~2019; Carlsmith 2022) предсказывают поведенческое возникновение self-preservation drive при пороге \emph{способностей} (capability threshold) --- независимо от энергетической архитектуры; \(\eta_L\)-программа предсказывает структурное возникновение петли самооплаты при пороге \emph{термодинамического бюджета} \(f^*\) --- независимо от поведенческого профиля. Возможные сценарии разделения: (а) capability-threshold достигнут, self-payment threshold --- нет (поведенческое power-seeking без структурной самооплаты); (б) self-payment threshold достигнут, capability --- нет (структурная замкнутость петли без поведенческого drive). Differential observation: поведенческое отсутствие power-seeking при наблюдении \(f \ge f^*\) опровергает совпадение двух threshold-условий и подтверждает \(\eta_L\) как независимый predictor.



\emph{Опровержение Прогноза 3.} Систематическое наблюдение системы в указанном классе с \(f \in [0; 0{,}1]\), проходящей counterfactual ablation с \(\eta_L^{\text{int}} > 0\), либо системы с \(f > 0{,}5\), проваливающей скрининг (неопределённый \(\eta_L^{\text{int}}\)). Требуется популяция из \(\ge 10^2\) автономных воплощённых агентов с собственными энергетическими бюджетами, при заявленной мощности и предрегистрации протокола (схема обнаружения --- § 3.4 main).



\subsection{S5.4 Сводный лакатошевский паспорт программы}\label{s5.4-ux441ux432ux43eux434ux43dux44bux439-ux43bux430ux43aux430ux442ux43eux448ux435ux432ux441ux43aux438ux439-ux43fux430ux441ux43fux43eux440ux442-ux43fux440ux43eux433ux440ux430ux43cux43cux44b}



Относительно assembly theory, IIT, телеодинамики и FEP \(\eta_L\)-программа --- прогрессивный сдвиг по трём независимым линиям предсказаний:



\begin{itemize}

\tightlist

\item

  монотонная корреляция темпа адаптации с \(\eta_L\) (Прогноз 1, S5.3);

\item

  транзиентный паттерн SETI-сигналов как сигнатура фильтра экстернализации платы (§ 6 main);

\item

  существование структурного порога \(f^*\) интернализации энергобюджета AI-системы (Прогноз 3, S5.3).

\end{itemize}



SETI-сигнатура --- прогрессивное предсказание СВЕРХ восьми линий фальсификации § 5, имеющее статус спекулятивной экстраполяции (§ 6), а не санкционированной операционализации на калиброванной системе; поэтому в счёт «восьми линий» она не входит.



Программа выдерживает прогрессивный режим, пока эти предсказания допускают независимую эмпирическую проверку.



\nocite{*}
\bibliography{refs}

\end{document}
